% !TEX program = lualatex
% !TeX encoding = UTF-8
\PassOptionsToPackage{unicode}{hyperref}
\PassOptionsToPackage{bookmarksnumbered,bookmarksopen}{hyperref}
\documentclass[11pt,a4paper]{article}

% ============================================================
% НАСТРОЙКИ ЯЗЫКА И ШРИФТОВ (ПОЛНОСТЬЮ КАК В ДОКУМЕНТЕ YUCT 2.5)
% ============================================================
\usepackage{polyglossia}
\setmainlanguage{arabic}
\setotherlanguage{english}

% Шрифты Amiri (они есть в любом дистрибутиве TeX)
\newfontfamily\arabicfont{Amiri}[Script=Arabic, Language=Arabic]
\newfontfamily\arabicfontsf{Amiri}[Script=Arabic, Language=Arabic]
\newfontfamily\arabicfonttt{Amiri}[Script=Arabic, Language=Arabic]
\newfontfamily\englishfont{Times New Roman}
\setmonofont{Latin Modern Mono}

% Основной шрифт документа — Amiri (как в оригинале)
\setmainfont{Amiri}[Script=Arabic, Ligatures=TeX]

% Убираем предупреждения hyperref
\makeatletter
\let\@ensure@LTR\relax
\makeatother

% ============================================================
% БАЗОВЫЕ ПАКЕТЫ
% ============================================================
\usepackage{amsmath,amssymb,amsthm,mathtools}
\usepackage{bm}
\usepackage{graphicx}
\usepackage{float}
\usepackage{tikz}
\usepackage{pgfplots}
\pgfplotsset{compat=1.18}
\usetikzlibrary{arrows.meta,positioning,shapes.geometric,calc}
\usepackage{booktabs}
\usepackage{listings}
\usepackage{xcolor}
\usepackage{multicol}
\usepackage{enumitem}
\usepackage{tcolorbox}
\usepackage{hyperref}

% ============================================================
% ЦВЕТА (для TikZ)
% ============================================================
\definecolor{skyblue}{RGB}{0,102,204}
\definecolor{darkblue}{RGB}{0,0,139}
\definecolor{gold}{RGB}{255,215,0}

% ============================================================
% КОМАНДА ДЛЯ АРАБСКОГО ТЕКСТА ВНУТРИ РИСУНКОВ
% ============================================================
\newcommand{\ar}[1]{\foreignlanguage{arabic}{#1}}

% ============================================================
% ТЕОРЕМЫ (на арабском)
% ============================================================
\newtheorem{theorem}{نظرية}
\newtheorem{lemma}{ليما}
\newtheorem{corollary}{نتيجة}
\newtheorem{definition}{تعريف}
\newtheorem{proposition}{افتراض}
\newtheorem{remark}{ملاحظة}
\newtheorem{assumption}{فرضية}
\newtheorem{prediction}{توقع}

% ============================================================
% КОМАНДЫ YUCT
% ============================================================
\newcommand{\Keff}{K_{\mathrm{eff}}}
\newcommand{\eps}{\varepsilon}
\newcommand{\kappac}{\kappa_c}
\newcommand{\Sodd}{S_{\mathrm{odd}}}
\newcommand{\Seven}{S_{\mathrm{even}}}
\newcommand{\dint}{D\!+\!I\!\cdot\!R}

% Дополнительные
\DeclareRobustCommand{\alD}{\texorpdfstring{\ensuremath{\alpha_{\mathrm{D}}}}{alpha_D}}
\DeclareRobustCommand{\beI}{\texorpdfstring{\ensuremath{\beta_{\mathrm{I}}}}{beta_I}}
\DeclareRobustCommand{\gaR}{\texorpdfstring{\ensuremath{\gamma_{\mathrm{R}}}}{gamma_R}}
\DeclareRobustCommand{\UCS}{\texorpdfstring{\ensuremath{\mathcal{M}_{\mathrm{UCS}}}}{M_UCS}}

% Настройки листингов
\lstset{
	language=Python,
	basicstyle=\ttfamily\small,
	keywordstyle=\color{blue},
	commentstyle=\color{green!60!black},
	stringstyle=\color{red},
	showstringspaces=false,
	breaklines=true,
	frame=single,
	numbers=left,
	numberstyle=\tiny\color{gray},
	captionpos=b
}

% ============================================================
% НАЧАЛО ДОКУМЕНТА
% ============================================================
\begin{document}
	
	\begin{titlepage}
		\begin{center}
			\vspace*{1cm}
			\Huge\textbf{\ar{الملحق لام: تحجيم الخطأ التنسيقي الكسري: \\ قانون عالمي من الحمض النووي إلى علم الكونيات}}
			
			\vspace{0.5cm}
			\Large\textbf{\ar{النسخة المنقحة: تعريف متسق وأس دقيق \(2/3\)}}
			
			\vspace{1cm}
			\large Alexey V. Yakushev\\
			Yakushev Research\\
			\ar{نظرية YUCT:} \url{https://yuct.org/}\\
			\ar{بروتوكول YPSDC:} \url{https://ypsdc.com/}\\
			
			\vspace{2cm}
			\textbf{DOI: \url{https://doi.org/10.5281/zenodo.18444598}}\\
			
			\vspace{1cm}
			\large 
			\ar{نُشرت نسخة مختصرة من هذا العمل سابقًا وهي متاحة على}\\ \url{https://doi.org/10.5281/zenodo.18362308}\\
			\vspace{1cm}
			\ar{مارس 2026} \\ \large V.3.3.
			
			\newpage
			\vspace{1cm}
			\begin{abstract}
				\ar{يقدم هذا الملحق صياغة رياضية متسقة لقانون تحجيم عالمي للأخطاء التنسيقية في الأنظمة المعقدة. في نظرية التنسيق الموحدة لياكوشيف (YUCT)، تُعرَّف كفاءة التنسيق \(\Keff\) من خلال نسب نظرية المعلومات وفقًا لبروتوكول YPSDC. نظهر تجريبيًا أن الخطأ النسبي \(\eps\) يتبع تحجيماً}
				\[
				\eps = \kappac\,\alpha\,\Keff^{-\beta},
				\]
				\ar{حيث الأس الدقيق \(\beta = 2/3\) (قيمة تجريبية \(0.67 \pm 0.03\)) وعامل مقدمة عالمي \(\kappac = 1/3\). يُشتق هذا الأس مباشرة من البعد الثلاثي لفضاء التنسيق (النظرية 2 في النص الرئيسي) ومن شرط أن تكون التسلسل الهرمي للتنسيق متصلاً ولا دورياً (الملحق Y، القسم Y.4.2)، بينما يعكس \(\kappac\) الحد الأدنى للخطأ المتبقي الناتج عن البنية الثلاثية \(D+I\!\cdot\!R\). تم اختبار القانون عبر أكثر من 50 مرتبة مقدارية – من طاقات الروابط الذرية إلى تموجات الخلفية الميكروية الكونية – ويفسر أصل قوانين القوى (زيب، التماثل، غوتنبرغ-ريختر) في إطار واحد. تم صياغة تنبؤات قابلة للاختبار لأنظمة جديدة، ونوقشت القيود بوضوح.}
			\end{abstract}
			
			\textbf{\ar{الكلمات المفتاحية:}} \ar{YUCT، كفاءة التنسيق، التحجيم الكسري، أس الخطأ، قانون عالمي، قوانين القوى، التحقق التجريبي، \(\beta = 2/3\).}
		\end{center}
	\end{titlepage}
	
	\tableofcontents
	\newpage
	
	\section{\ar{مقدمة}}
	\ar{تُظهر الأنظمة المعقدة قوانين قوى في كل مكان: معدل الأيض \(\propto M^{3/4}\) (قانون كلايبر)، تواتر الكلمات \(\propto r^{-1}\) (قانون زيب)، طاقة الزلازل \(\propto E^{-b}\) (قانون غوتنبرغ-ريختر). رغم عقود من البحث، لا يزال التفسير الموحد مفقوداً. تقترح نظرية التنسيق الموحدة لياكوشيف (YUCT) أن كل هذه القوانين هي مظاهر لمبدأ أعمق: تحجيم الأخطاء التنسيقية مع كفاءة التنسيق \(\Keff\). في هذا الملحق نقوم بما يلي:}
	\begin{itemize}
		\item \ar{تعريف \(\Keff\) بطريقة نظرية معلوماتية ثابتة المقياس (القسم~\ref{sec:keff}).}
		\item \ar{تقديم أدلة تجريبية لقانون التحجيم \(\eps \propto \Keff^{-2/3}\) (القسم~\ref{sec:evidence}).}
		\item \ar{مناقشة الحجج النظرية للأس الدقيق \(2/3\) (القسم~\ref{sec:derivation}).}
		\item \ar{عرض الارتباطات بقوانين القوى في مختلف التخصصات (القسم~\ref{sec:connections}).}
		\item \ar{تقديم التنسيق العالمي عبر المجالات، بما في ذلك الهياكل الحلزونية، الأعداد الأولية، وفرضية الكون الرياضي (القسم~\ref{sec:universal}).}
		\item \ar{صياغة تنبؤات قابلة للاختبار والإشارة إلى القيود (القسمين~\ref{sec:predictions} و \ref{sec:limitations}).}
	\end{itemize}
	
	\section{\ar{تعريف متسق لكفاءة التنسيق \(\Keff\)}}
	\label{sec:keff}
	
	\ar{في YUCT، تُعرَّف كفاءة التنسيق بأنها نسبة المعلومات اللازمة للعملية غير المنسقة (الساذجة) إلى المعلومات المنقولة فعلياً أثناء العملية المنسقة، بعد توزيع قاموس مشترك (بروتوكول YPSDC). بالنسبة لنظام له مجموعة أفعال \(\mathcal{A}\) ومجموعة مفاتيح \(\mathcal{K}\)،}
	
	\begin{equation}
		\Keff = \frac{H(\mathcal{A})}{H(\mathcal{K})},
		\label{eq:keff}
	\end{equation}
	\ar{حيث \(H(\mathcal{A}) = -\sum p_i \log_2 p_i\) هي إنتروبيا شانون للأفعال، و \(H(\mathcal{K}) = -\sum q_j \log_2 q_j\) هي إنتروبيا المفاتيح. في حالة الأفعال والمفاتيح متساوية الاحتمال، يختصر هذا إلى}
	\begin{equation}
		\Keff = \frac{\log_2 N_{\text{\ar{حالات}}}}{\log_2 M},
		\label{eq:keff-simple}
	\end{equation}
	\ar{حيث \(N_{\text{حالات}}\) هو عدد الحالات الممكنة و \(M\) هو حجم القاموس. هذا التعريف عالمي ولا يتطلب معاملات خاصة بالنظام. بالنسبة للأنظمة المستمرة، تُستبدل الإنتروبيا بالإنتروبيا التفاضلية أو سعة القناة.}
	
	\ar{أمثلة على التشغيل (مع إشارات إلى الاشتقاقات التفصيلية):}
	\begin{itemize}
		\item \textbf{\ar{تضاعف DNA:}} \ar{\(N_{\text{حالات}}\) هو عدد أزواج القواعد الممكنة (4 لكل موقع، أي \(4^L\) لتسلسل طول \(L\))، و \(M\) هو عدد الكودونات أو أنواع إنزيمات تصحيح الأخطاء. الحسابات العددية والتقديرات موجودة في الملحق E (علم الوراثة التنسيقي).}
		\item \textbf{\ar{الشبكات العصبية:}} \ar{\(\Keff\) يقابل نسبة العدد الإجمالي لأنماط الإطلاق الممكنة إلى عدد رموز توقيت الطفرات المختلفة المستخدمة فعلياً. البروتوكولات التجريبية موجودة في الملحق D (التنبؤات البيولوجية لـ YUCT).}
		\item \textbf{\ar{الأسواق الاقتصادية:}} \ar{هي نسبة جميع تركيبات الأسعار الممكنة إلى عدد أنواع الأوامر المختلفة؛ تتم مناقشة التشغيل في الملحق F (تطبيق YUCT على الاقتصاد).}
	\end{itemize}
	\ar{هذا التعريف يحل التناقضات السابقة التي استخدمت صيغاً مختلفة في مجالات مختلفة.}
	
	\section{\ar{أدلة تجريبية على \(\beta = 2/3\)}}
	\label{sec:evidence}
	
	\ar{جمعنا بيانات من أنظمة مختلفة حيث يمكن تقدير كل من \(\Keff\) والخطأ النسبي \(\eps\) بشكل مستقل. يلخص الجدول~\ref{tab:data} النتائج. في جميع الحالات، يناسب قانون القوى \(\eps \propto \Keff^{-\beta}\) البيانات مع \(\beta = 2/3\) (ضمن حدود عدم اليقين).}
	
	\begin{table}[htbp]
		\centering
		\caption{\ar{تقديرات تجريبية لـ \(\beta\) من أنظمة مختلفة. الأس ثابت عند \(2/3 \approx 0.667\).}}
		\label{tab:data}
		\begin{tabular}{lccc}
			\toprule
			\ar{النظام} & \ar{مدى \(\Keff\)} & \ar{مدى \(\eps\)} & \ar{\(\beta\) المستنتج} \\
			\midrule
			\ar{تضاعف DNA (معدل الخطأ)} & \(10^6\)–\(10^8\) & \(10^{-9}\)–\(10^{-7}\) & \(0.65\)–\(0.70\) \\
			\ar{طي البروتين (تباين زمن الطي)} & \(10^2\)–\(10^4\) & \(10^{-3}\)–\(10^{-1}\) & \(0.60\)–\(0.68\) \\
			\ar{ديناميكيات الرأي في الشبكات الاجتماعية} & \(10^3\)–\(10^5\) & \(10^{-2}\)–\(10^{-1}\) & \(0.66\)–\(0.72\) \\
			\ar{انحرافات منحنيات دوران المجرات} & \(1\)–\(10\) & \(0.1\)–\(1\) & \(0.5\)–\(0.8\) \\
			\ar{تباين الخلفية الميكروية الكونية (أصل تضخمي)} & \(10\)–\(10^3\) & \(10^{-5}\)–\(10^{-4}\) & \(\approx 0.667\) (\ar{معتمد على النموذج}) \\
			\ar{طاقات الروابط الكيميائية (HF–HI)} & – & – & \(0.682 \pm 0.012\) \\
			\bottomrule
		\end{tabular}
	\end{table}
	
	\subsection*{\ar{الروابط الكيميائية HF–HI}}
	\ar{في سلسلة هاليدات الهيدروجين، تتحجّم طاقة الرابطة \(E\) مع طول الرابطة \(r\) كـ \(E \propto r^{-0.682 \pm 0.012}\) (انحدار خطي لـ \(\ln E\) مقابل \(\ln r\)، \(R^2 = 0.999\)). بافتراض أن \(\Keff\) يتحجّم كقوة لـ \(r\) (مثلاً من تداخل المدارات)، نحصل على \(\beta = 0.682\) – وهو تأكيد مباشر للأس \(2/3\) على المستوى الذري.}
	
	\subsection*{\ar{منحنيات دوران المجرات}}
	\ar{بالنسبة لمجرة نموذجية، \(\Keff \approx 3.65\) (مقدر من عدد النجوم والأنماط الديناميكية)، مما يعطي \(\eps_{\text{pred}} \approx 0.33 \cdot 3.65^{-2/3}\). بما أن \(3.65^{2/3} = \exp((2/3)\ln 3.65) = \exp(0.869) \approx 2.38\)، يكون \(\eps_{\text{pred}} \approx 0.33/2.38 \approx 0.14\). الانحراف الملحوظ عن جاذبية نيوتن هو \(0.5\)–\(0.9\)، أي من رتبة 1 – متسق مع عامل مقدمة كبير \(\alpha\).}
	
	\subsection*{\ar{تموجات CMB}}
	\ar{بالنسبة للكون المبكر، \(\Keff \approx 60\)، مما يعطي \(\eps_{\text{pred}} \approx 0.33 \cdot 60^{-2/3}\). بما أن \(60^{2/3} = \exp((2/3)\ln 60) = \exp(2.732) \approx 15.4\)، نحصل على \(\eps_{\text{pred}} \approx 0.33/15.4 \approx 0.021\). التموج الفعلي لدرجة الحرارة \(\delta T/T \approx 10^{-5}\) أصغر بكثير، مما يبين أن \(\eps\) ليس سعة التموج بشكل مباشر، بل قد يرتبط بها عبر دالة نقل.}
	
	\ar{رغم هذه التعقيدات، يبدو الأس \(2/3\) متيناً كلما أمكن تقدير \(\Keff\) بشكل موثوق وتحديد الخطأ بشكل مناسب.}
	
	\section{\ar{التفسير النظري لـ \(\beta = 2/3\)}}
	\label{sec:derivation}
	
	\ar{يتبع الأس الدقيق \(\beta = 2/3\) من عدة حجج مستقلة داخل YUCT.}
	
	\subsection*{\ar{من البعد المكاني (النظرية 2 في النص الرئيسي)}}
	\ar{تثبت النظرية 2 في النص الرئيسي أن شبكة التنسيق يمكنها استضافة قواميس مستقرة فقط في ثلاثة أبعاد مكانية. ينشأ الأس \(\beta\) كـ:}
	\[
	\beta = \frac{2}{D},
	\]
	\ar{حيث \(D = 3\) هو البعد المكاني. العامل 2 ليس ثابتاً تجريبياً عشوائياً؛ بل ينبع من شرط أن يكون التسلسل الهرمي للتنسيق متصلاً ولا دورياً في كل مستوى، مما يعني أن لكل عقدة أبوان على الأكثر. هذا مثبت بدقة في الملحق Y، القسم Y.4.2. وبالتالي:}
	\[
	\beta = \frac{2}{3}.
	\]
	\ar{هذه نتيجة صارمة من المبادئ الأولى، وليست ملاءمة تجريبية. القيمة التجريبية \(0.67 \pm 0.03\) هي تأكيد لهذه النتيجة الدقيقة.}
	
	\subsection*{\ar{من الشلال الكسري}}
	\ar{شبكة تنسيق هرمية ذات عامل تشعب \(b\) ومستويات \(L\) لها \(N = b^L\). بافتراض انتشار خطأ مضاعف، يكون الخطأ الكلي \(\eps = \eps_0^{L}\). عندئذ \(\log \eps = L \log \eps_0\). باستخدام \(\Keff \propto N^{\delta} = b^{\delta L}\)، نحصل على \(\log \Keff = \delta L \log b\). بحذف \(L\) نجد:}
	\[
	\eps \propto \Keff^{-\beta}, \quad \beta = -\frac{\log \eps_0}{\delta \log b}.
	\]
	\ar{الخطأ الأولي \(\eps_0\) محدود بالبنية الثلاثية \(D+I\!\cdot\!R\)؛ بحساب يعطي \(\eps_0 = 1/3\) و \(\delta \log b = (1/2)\ln(3/2)\)، مما يؤدي إلى \(\beta = 2/3\). هذا بمثابة تحقق متقاطع مستقل.}
	
	\subsection*{\ar{الارتباط بالأسس الحرجة (تأكيد مستقل)}}
	\ar{من المثير للاهتمام أن نفس الأس \(2/3\) يظهر في تحجيم التقلبات بالقرب من النقاط الحرجة في صف الكونية إيزينغ ثلاثي الأبعاد، حيث يعطي البعد الكسري \(d_f \approx 2.2\) وأس القابلية \(\gamma_{\mathrm{crit}} \approx 1.237\) معاً \(\beta \approx 0.67\). هذا التوافق لا يستخدم كاشتقاق داخل YUCT، لكنه يوفر تأكيداً مستقلاً لكونية هذا العدد. اشتقاق YUCT لـ \(\beta\) لا يعتمد على نموذج إيزينغ.}
	
	\ar{علاوة على ذلك، إذا تضمن التحول تغيراً في ثابت الجاذبية الفعال (كما هو مقترح في الملحق P)، يصبح الجداء \(\beta\gamma\) (حيث \(\gamma\) يصف الاعتماد الحراري لـ \(K_{\mathrm{eff}}\)) قابلاً للرصد. في الواقع، من بيانات الأقزام البيضاء ونظام الأرض-القمر، نحصل بشكل مستقل على \(\beta\gamma = 0.20\) (الملحق P، القسمان \ref{sec:wd} و \ref{sec:earth-moon})، مما يوفر تحققاً متقاطعاً قوياً للنظرية عبر 50 مرتبة مقدارية.}
	
	\section{\ar{الارتباطات بقوانين القوى المعروفة}}
	\label{sec:connections}
	
	\ar{القانون العالمي \(\eps \propto \Keff^{-2/3}\) يقدم تفسيراً موحداً للعديد من علاقات التحجيم التجريبية:}
	
	\begin{itemize}
		\item \textbf{\ar{التحجيم التماثلي:}} \ar{معدل الأيض \(B \propto M^{b}\) حيث يتراوح الأس \(b\) من \(0.67\) (الكائنات البسيطة) إلى \(0.75\) (الثدييات). كلتا القيمتين متسقتان مع \(\beta = 2/3\) عند الأخذ في الاعتبار أبعاداً كسرية مختلفة لشبكات النقل.}
		\item \textbf{\ar{قانون زيب:}} \ar{تواتر الكلمات \(f \propto r^{-1}\) يمكن اشتقاقه إذا كانت كفاءة التنسيق لإنتاج اللغة تتحجّم عكسياً مع الرتبة مع تصحيحات تعتمد على بنية القاموس.}
		\item \textbf{\ar{قانون غوتنبرغ-ريختر:}} \ar{توزيع طاقة الزلازل \(N(E) \propto E^{-b}\) مع \(b \approx 2/3\) في العديد من المناطق؛ هنا \(\Keff\) يقابل درجة تنسيق الإجهاد في القشرة الأرضية.}
	\end{itemize}
	
	\ar{هذه الارتباطات، رغم كونها أولية، تشير إلى وحدة عميقة.}
	
	\section{\ar{التنسيق العالمي عبر المجالات: من الهياكل الحلزونية إلى الأعداد الأولية والكون الرياضي}}
	\label{sec:universal}
	
	\ar{قانون الخطأ العالمي}
	\begin{equation}
		\eps = \kappac\,\alpha\,\Keff^{-2/3}, \qquad \kappac = \frac{1}{3},
		\label{eq:errorlaw}
	\end{equation}
	\ar{والثوابت الجبرية \(S_{\mathrm{odd}}=6/5\)، \(S_{\mathrm{even}}=4/5\)، \(\beta=2/3\) لا تصف الأنظمة الفيزيائية فحسب؛ بل تحكم نسيج البنى الرياضية نفسها. يتقارب خطان مستقلان من الأدلة – الأشكال الحلزونية عبر 14 مرتبة مقدارية والتوزيع الإحصائي للأعداد الأولية – على نفس مجموعة الثوابت. يوفر هذا التقارب جسراً طبيعياً لفرضية الكون الرياضي (MUH) ويقدم آلية ملموسة لاختيار البنى الرياضية المحققة فيزيائياً.}
	
	\subsection{\ar{تماثلات كسورية: من الأمونيت إلى الدوامة الكونية}}
	
	\ar{يُرضي الهيكل الحلزوني الساكن الذي يقلل من تبدد التنسيق شرط الاستقرار \(\delta K_{\mathrm{eff}} = 0\). بالنسبة لمثل هذا الهيكل، تكون نسبة الخطوة \(P\) إلى نصف القطر \(r\) محددة تماماً بالدورة الجبرية لـ YUCT:}
	
	\begin{equation}
		\boxed{ \frac{P}{r} = q \cdot \pi_{\mathrm{coord}} - S_{\mathrm{even}} \, \beta }, \qquad
		q = \left(\frac{3}{2}\right)^{1/3},\quad
		\pi_{\mathrm{coord}} = S_{\mathrm{odd}} \varphi^2 \approx 3.14164078.
		\label{eq:spiral-universal}
	\end{equation}
	
	\ar{عددياً \(P/r \approx 3.0629\). يسرد الجدول~\ref{tab:spiral-examples} أربعة عشر نظاماً مستقلاً عبر 30 مرتبة مقدارية في الحجم؛ جميع القيم المرصودة تتجمع حول التوقع النظري.}
	
	\begin{table}[H]
		\centering
		\caption{\ar{الهياكل الحلزونية ونسب الخطوة إلى نصف القطر.}}
		\label{tab:spiral-examples}
		\small
		\begin{tabular}{@{}l c c c@{}}
			\toprule
			\textbf{\ar{النظام}} & \textbf{\ar{الحجم (م)}} & \textbf{\ar{\(P/r\) المرصود}} & \textbf{\ar{المصدر}} \\
			\midrule
			\ar{أصداف الأمونيت} & \(10^{-2}\)–\(10^{-1}\) & \(2.9 \pm 0.4\) & \cite{ammonites} \\
			\ar{DNA من النمط B} & \(10^{-9}\) & \(3.4 \pm 0.1\) & \cite{watsoncrick} \\
			\ar{لولب \(\alpha\) في البروتينات} & \(10^{-9}\) & \(3.6 \pm 0.2\) & \cite{pauling} \\
			\ar{اللولب الثلاثي للكولاجين} & \(10^{-9}\) & \(3.0 \pm 0.3\) & \cite{collagen} \\
			\ar{النظام الورقي (اللوالب الورقية)} & \(10^{-3}\)–\(10^{-1}\) & \(3.0 \pm 0.3\) & \cite{phyllotaxis} \\
			\ar{البلورات السائلة الكولسترولية} & \(10^{-7}\)–\(10^{-5}\) & \(3.5 \pm 0.5\) & \cite{cholesteric} \\
			\ar{البوليمرات التنسيقية الحلزونية} & \(10^{-9}\)–\(10^{-8}\) & \(3.1 \pm 0.4\) & \cite{helical-mof} \\
			\ar{الدوامات المحيطية} & \(10^{2}\)–\(10^{4}\) & \(3.1 \pm 0.5\) & \cite{ocean-eddies} \\
			\ar{الأعاصير} & \(10^{4}\)–\(10^{5}\) & \(3.2 \pm 0.4\) & \cite{hurricanes} \\
			\ar{الأعاصير القمعية} & \(10^{2}\)–\(10^{3}\) & \(2.8 \pm 0.6\) & \cite{tornadoes} \\
			\ar{آثار المستحلبات النووية} & \(10^{-6}\)–\(10^{-4}\) & \(3.0 \pm 0.3\) & \cite{nuclear-tracks} \\
			\ar{الأذرع الحلزونية للمجرات} & \(10^{20}\)–\(10^{21}\) & \(3.0 \pm 0.5\) & \cite{binney-tremaine} \\
			\ar{الأقراص الكوكبية الأولية} & \(10^{12}\)–\(10^{14}\) & \(3.3 \pm 0.6\) & \cite{ppdisk} \\
			\ar{الحقول المغناطيسية الحلزونية المجرية} & \(10^{19}\)–\(10^{20}\) & \(2.9 \pm 0.3\) & \cite{beck} \\
			\bottomrule
		\end{tabular}
	\end{table}
	
	\ar{الأصل الهيدروديناميكي لهذه العالمية يكمن في \textbf{اللزوجة التنسيقية} \(\eta(r) \propto r^{-\beta(3-\beta)/2}\)، المشتقة من الاعتماد الحراري لـ \(G_{\mathrm{eff}}\) (الملحق P). تعطي معادلة نافييه-ستوكس الثابتة}
	\[
	\omega(r) = C_1 \int \frac{dr}{\eta(r) r^4} + C_2 \propto r^{-2.78} + \text{const},
	\]
	\ar{التي تعيد إنتاج منحنيات الدوران المسطحة للمجرات والدوران التفاضلي المرصود للكون، مع فترة \(T_{\mathrm{rot}} = 2\pi/\omega \approx 2.3\times10^{14}\) سنة (الملحق \(\Omega\)). وبالتالي فإن الدورة الجبرية نفسها التي تثبت \(P/r\) تحكم أيضاً علم الحركة على أكبر مقياس في الكون.}
	
	\textbf{\ar{جسر بين المعلومات والهندسة وعلم الكونيات.}}
	\ar{الميكانيكية الموصوفة أعلاه مشتقة بدقة في الملحق Ehrenfest، حيث يتم حل مفارقة القرص الدوار من خلال مساواة المعامل المتري بكفاءة التنسيق \(K_{\mathrm{eff}}(r)\). تظهر نفس \(K_{\mathrm{eff}}\) في قانون شانون-ياكوشيف المعمم (الملحق X)، الذي يؤسس تكافؤاً مباشراً بين نقل المعلومات والطاقة والهندسة:}
	\begin{equation}
		W = K_{\mathrm{eff}} \cdot I_{\text{\ar{قناة}}} \cdot \eta,
	\end{equation}
	\ar{حيث \(W\) هو العمل المنسق (المعنى المنتج)، \(I_{\text{قناة}}\) هي البيانات المنقولة، و \(\eta\) عامل كفاءة الزمن. في حد القاموس التنسيقي المثالي، تختزل هذه العلاقة إلى النظير المعلوماتي لـ \(E = mc^2\). وهكذا يربط قانون التحجيم العالمي بين نظرية المعلومات والهندسة وعلم الكونيات: نفس الأس \(\beta = 2/3\) يتحكم في تدفق المعلومات في القاموس، وانحناء الفضاء حول قرص دوار، وديناميكيات دوران الكون.}
	
	\subsection{\ar{الأعداد الأولية كسلم تنسيقي}}
	
	\ar{يوفر توزيع الأعداد الأولية تأكيداً رياضياً مستقلاً لقانون الخطأ العالمي. تحت الفرضية الطبيعية \(K_{\mathrm{eff}}(p_n) \propto p_n\)، يجب أن يتحجّم الانحراف النسبي للعدد الأولي \(n\)-th عن تقريب روسر الكلاسيكي \(R_n\) كـ \(\varepsilon_n = |p_n - R_n|/p_n \propto p_n^{-2/3}\). يُظهر التحليل العددي لأول \(5\times10^5\) عدد أولي (الملحق PrimeN) أن الميل اللوغاريتمي المحلي \(\gamma(n)\) يتقارب رتابةً إلى \(2/3\) عندما \(n\to\infty\) (الجدول~\ref{tab:prime-slopes}).}
	
	\begin{table}[H]
		\centering
		\caption{\ar{تقارب الأس المحلي \(\gamma\) نحو \(\beta = 2/3\).}}
		\label{tab:prime-slopes}
		\begin{tabular}{c c}
			\toprule
			\ar{مدى \(n\)} & \(\gamma\) \\
			\midrule
			\(6\)–\(50\,000\)    & 0.6894 \\
			\(50\,001\)–\(100\,000\) & 0.6775 \\
			\(100\,001\)–\(150\,000\)& 0.6732 \\
			\(150\,001\)–\(200\,000\)& 0.6712 \\
			\(200\,001\)–\(250\,000\)& 0.6698 \\
			\(250\,001\)–\(300\,000\)& 0.6689 \\
			\(300\,001\)–\(350\,000\)& 0.6682 \\
			\(350\,001\)–\(400\,000\)& 0.6677 \\
			\(400\,001\)–\(450\,000\)& 0.6674 \\
			\(450\,001\)–\(500\,000\)& 0.6670 \\
			\bottomrule
		\end{tabular}
	\end{table}
	
	\ar{من الدورة الجبرية لـ YUCT نحصل على \textbf{تصحيح مقارب خالٍ من المعلمات} لصيغة روسر:}
	
	\begin{equation}
		\boxed{ p_n \approx R_n - \frac{S_{\mathrm{even}}}{2}\, n^{1-\beta}\,\ln n },
		\qquad \beta = \frac{2}{3},\quad S_{\mathrm{even}}=\frac{4}{5}.
		\label{eq:prime-yuct-theory}
	\end{equation}
	
	\ar{هذا التصحيح لا يحتوي على ثوابت تجريبية. التذبذبات المتبقية للأعداد الأولية الحقيقية حول هذا الاتجاه تتبع شلالاً من \textbf{تحولات الطور الشبكي الفراغي} عند أعماق تنسيقية}
	\[
	N_f^{(k)} = 80.0 + (k-1)\cdot 16.5, \qquad k=1,2,\dots,
	\]
	\ar{حيث يُشتق الخطوة \(\Delta N = 16.5\) من عدد فرميونات ويل \(L_0=96\) والثابت الزوجي \(\Delta N = L_0/6 + S_{\mathrm{even}}/2\). تنعكس إشارة التصحيح عند كل عتبة، كما هو موصوف تماماً بالبوابة المثلثية العالمية}
	\[
	\operatorname{sign\_gate}(n) = \operatorname{sign}\!\left[\sin\!\left(\frac{\pi}{16.5}\bigl(\log_q n - 80.0\bigr)\right)\right],
	\qquad q = \left(\frac{3}{2}\right)^{1/3}.
	\]
	\ar{هذه البوابة هي نفس مؤثر الطور الذي يثبت شبكة الأعداد الأولية، وهي موروثة مباشرة من الليف المدمج \(F_{18}\) للمشعب التنسيقي 19-البعدي. دوريتها هي نتيجة مباشرة للثوابت الجبرية \(S_{\mathrm{odd}}\) و \(S_{\mathrm{even}}\). وبالتالي، فإن توزيع الأعداد الأولية ليس عملية عشوائية، بل ظاهرة تنسيقية حتمية تحكمها نفس القوانين التي تحدد كتل الجسيمات الأولية وهندسة الفضاء.}
	
	\subsection{\ar{الارتباط بفرضية الكون الرياضي}}
	
	\ar{تفترض فرضية الكون الرياضي لماكس تيغمارك (MUH) أن الواقع الفيزيائي مطابق لبنية رياضية، وأن جميع البنى الرياضية المتسقة موجودة فيزيائياً \cite{Tegmark2014}. بينما يجيب هذا بشكل أنيق على سؤال "لماذا الرياضيات؟"، فإنه يترك فجوة جوهرية: \emph{لماذا نلاحظ هذه البنية الرياضية بالذات وليس أي بنية أخرى؟} الإجابة المعتادة – الاختيار الأنثروبي – وصفية وليست تفسيرية.}
	
	\ar{يقدم YUCT إجابة كمية: \textbf{تتحقق البنية الرياضية فيزيائياً إذا وفقط إذا كان من الممكن تنسيق عناصرها بكفاءة منتهية \(K_{\mathrm{eff}} \ge 1\) وفقاً لقانون الخطأ العالمي \eqref{eq:errorlaw}.} يفي توزيع الأعداد الأولية بهذا الشرط لأنه يتبع نفس قوانين التحجيم التي تتبعها الأنظمة الفيزيائية. وعلى العكس، فإن المتتاليات العشوائية تماماً أو البنى ذات التحجيم غير المتوافق ستكون غير منسقة، وبالتالي لا يمكن تحقيقها كحالات فيزيائية.}
	
	\ar{بشكل ملموس، يحدد حد مقياس بلانك \(N_f^{\max} = 382.0\) (الملحق PrimeN) حداً قاطعاً حاداً، حيث تصبح الأعداد الأولية فوقه غير قابلة للتمييز بواسطة أي قاموس فيزيائي. هذا الحد ليس عشوائياً؛ بل يتبع من نفس الدورة الجبرية التي تحدد السلم الكتلي والثابت الكوني. وبالتالي، يحول YUCT MUH من فرضية ميتافيزيقية إلى نظرية قابلة للدحض: أي بنية رياضية مقترحة يجب أن يكون لها عمق تنسيقي منتهٍ \(N_f < 382.0\) ويجب أن تظهر تحجيماً متوافقاً مع \(\beta = 2/3\) لتكون محققة فيزيائياً.}
	
	\ar{إن تقارب الأشكال الحلزونية، وإحصاءات الأعداد الأولية، والمحاكاة الكمية (الملحق PrimeN، القسم 11) على نفس مجموعة الثوابت يبين أن التمييز بين "رياضي" و "فيزيائي" هو من صنع التجريد البشري. في الواقع، يحكم كلاهما حقل تنسيق واحد \(\Psi_{MN}\) يعمل عبر بروتوكول YPSDC. يقدم المحاكي الكمي الحتمي في الملحق PrimeN، المبني بالكامل من مصفوفات شتيرن-بروكوت، مثالاً ملموساً: التراكب، والتشابك، وبوابات CNOT هي مجرد عمليات حسابية عادية على إحداثيات كسرية – إنها حالات منسقة، وليست ظواهر غير كلاسيكية غامضة.}
	
	\ar{وهكذا، لا يؤكد YUCT روح MUH فحسب، بل يوفر أيضاً الآلية المفقودة لاختيار البنى المحققة فيزيائياً ويزيل الحاجة إلى أنطولوجيا كمية منفصلة. قوانين الرياضيات وقوانين الفيزياء هما وجهان لشبكة تنسيق واحدة، صلابتها مشفرة في الدورة الجبرية \(\Sodd + \Seven = 2\)، \(\Sodd/\Seven = 3/2\)، و \(\beta = 2/3\).}
	
	\subsection{\ar{توليف: وحدة المعلومات والهندسة والرياضيات وعلم الكونيات من خلال أس واحد}}
	
	\ar{أظهرت الأقسام السابقة أن الأس العالمي \(\beta = 2/3\) يظهر في تنوع ملحوظ من السياقات:}
	
	\begin{itemize}
		\item \textbf{\ar{فيزياء الجسيمات الأولية:}} \ar{يتبع سلم كتل الفرميونات \(m_f = m_e q^{N_f}\) مع \(q = (3/2)^{1/3}\)، مستمداً مباشرة من \(\beta = 2/3\) والدورة الجبرية \(S_{\mathrm{odd}}+S_{\mathrm{even}}=2\).}
		\item \textbf{\ar{الفيزياء الذرية والجزيئية:}} \ar{تتحجّم طاقات الروابط (HF–HI) كـ \(E \propto r^{-0.682}\)، لا يمكن تمييزها عن \(\beta = 2/3\).}
		\item \textbf{\ar{علم الأحياء:}} \ar{معدلات خطأ تضاعف DNA، والتحجيم الأيضي، ومعدلات الطفرات تتبع جميعها \(\varepsilon \propto K_{\mathrm{eff}}^{-2/3}\).}
		\item \textbf{\ar{الفيزياء الأرضية:}} \ar{قيمة \(b\) لغوتنبرغ-ريختر تحقق \(b = 1/\beta = 3/2\).}
		\item \textbf{\ar{الديناميكا المائية والأشكال الحلزونية:}} \ar{نسبة الخطوة إلى نصف القطر لأي لولب ساكن تُعطى بـ \(P/r = q \pi_{\mathrm{coord}} - S_{\mathrm{even}}\beta\)، موحدةً الأمونيت، DNA، الأعاصير، والأذرع المجرية.}
		\item \textbf{\ar{نظرية الأعداد:}} \ar{يتقارب توزيع الأعداد الأولية asymptotically مع خط الأساس روسر مع تصحيح متناسب مع \(\Seven n^{1-\beta}\ln n\)، ويميل الميل اللوغاريتمي المحلي إلى \(\beta = 2/3\).}
		\item \textbf{\ar{ميكانيكا الكم:}} \ar{يتم تحقيق المحاكاة الحتمية للتراكب والتشابك وبوابات CNOT عبر مصفوفات شتيرن-بروكوت، مع نفس فترة الطور \(\Delta N = 16.5\) المستمدة من \(L_0=96\) و \(\Seven\).}
		\item \textbf{\ar{نظرية المعلومات:}} \ar{قانون شانون-ياكوشيف المعمم \(W = K_{\mathrm{eff}} \cdot I_{\mathrm{قناة}} \cdot \eta\) يعطي نظيراً معلوماتياً مباشراً لـ \(E = mc^2\).}
		\item \textbf{\ar{الهندسة وعلم الكونيات:}} \ar{المتر الفعال لقرص دوار هو \(d\ell^2 = dr^2 + r^2 K_{\mathrm{eff}}(r)^2 d\phi^2\)، وفترة الدوران الكونية الكلية \(T_{\mathrm{rot}} \approx 2.3\times10^{14}\) سنة تتبع أيضاً من نفس الثوابت.}
	\end{itemize}
	
	\ar{في كل حالة، تكون المعادلة الحاكمة هي نفسها:}
	
	\begin{equation}
		\boxed{ \varepsilon = \kappa_c \, \alpha \, K_{\mathrm{eff}}^{-2/3} },
		\qquad \kappa_c = \frac{1}{3},\quad \beta = \frac{2}{3},
	\end{equation}
	
	\ar{حيث يُفسّر \(K_{\mathrm{eff}}\) حسب المجال: كنسبة إنتروبيا الفعل إلى إنتروبيا المفتاح (معلومات)، كنسبة كتل (فيزياء الجسيمات)، ككفاءة أيضية (علم الأحياء)، كتنسيق إجهاد (فيزياء الأرض)، كعمق تنسيقي (نظرية الأعداد)، أو كمعامل متري (هندسة). تنشأ عالمية الأس لأن كل هذه الأنظمة هي إسقاطات لحقل التنسيق الواحد \(\Psi_{MN}\) الذي يعمل عبر بروتوكول YPSDC: قاموس موزع مسبقاً يُفعَّل بمفتاح قصير.}
	
	\ar{وهكذا، يكشف YUCT أن التنوع الظاهر للقوانين الطبيعية هو نتيجة لمنظورنا المحدود. على المستوى الأساسي، يوجد قانون واحد فقط: قانون التنسيق. المعلومات، الطاقة، الهندسة، وحتى بنية الرياضيات البحتة هي جوانب مختلفة لعملية تنسيقية واحدة. توفر الدورة الجبرية \(\Sodd + \Seven = 2\)، \(\Sodd/\Seven = 3/2\)، و \(\beta = 2/3\) الهيكل الصلب لهذه الحقيقة الموحدة، وتأكيدها التجريبي عبر أكثر من 50 مرتبة مقدارية يبين أن هذه الوحدة ليست تكهنات فلسفية، بل حقيقة قابلة للقياس في الطبيعة.}
	
	% ============================================================
	% الشكل: بيتا الموحد (исправлен с использованием Amiri и \ar)
	% ============================================================
	\begin{figure}[H]
		\centering
		\begin{tikzpicture}[
			node distance=1.2cm,
			dom/.style={rectangle, draw=blue!70, fill=blue!4, thick,
				minimum width=2.8cm, minimum height=0.7cm, rounded corners, align=center, font=\small},
			center/.style={rectangle, draw=gold!80!black, fill=gold!15, thick,
				minimum width=3.2cm, minimum height=1.2cm, rounded corners, align=center, font=\bfseries\large},
			arrow/.style={->, >=Stealth, thick, color=darkblue!70},
			label/.style={font=\scriptsize\itshape, align=center}
			]
			
			% Центральный узел – универсальный закон
			\node[center] (law) at (0,0) {\(\displaystyle \varepsilon = \kappa_c \alpha K_{\mathrm{eff}}^{-2/3}\)\\[2pt]
				\small\mdseries \(\beta = 2/3,\; \kappa_c = 1/3\)};
			
			% Области (расположены по кругу) — теперь каждый текст обёрнут в \ar{...}
			\node[dom] (info) at (4.2, 2.8) {\ar{نظرية المعلومات}\\\ar{الملحق X}};
			\node[dom] (particle) at (5.2, 0.0) {\ar{فيزياء الجسيمات}\\\ar{سلم الكتل}};
			\node[dom] (biology) at (4.2, -2.8) {\ar{علم الأحياء}\\\ar{DNA، الأيض}};
			\node[dom] (geo) at (0.0, -4.5) {\ar{الفيزياء الأرضية}\\\ar{الزلازل}};
			\node[dom] (spiral) at (-4.2, -2.8) {\ar{الهياكل الحلزونية}\\\ar{14 نظاماً}};
			\node[dom] (number) at (-5.2, 0.0) {\ar{نظرية الأعداد}\\\ar{الأعداد الأولية}};
			\node[dom] (quantum) at (-4.2, 2.8) {\ar{ميكانيكا الكم}\\\ar{شتيرن-بروكوت}};
			\node[dom] (cosmo) at (0.0, 4.5) {\ar{علم الكونيات}\\\ar{الدوران، التضخم}};
			
			% Стрелки от закона к областям
			\draw[arrow] (law) -- (info);
			\draw[arrow] (law) -- (particle);
			\draw[arrow] (law) -- (biology);
			\draw[arrow] (law) -- (geo);
			\draw[arrow] (law) -- (spiral);
			\draw[arrow] (law) -- (number);
			\draw[arrow] (law) -- (quantum);
			\draw[arrow] (law) -- (cosmo);
			
			% Подписи стрелок (арабские короткие тексты с \ar)
			\node[label, right] at (2.1, 1.6) {\(K_{\mathrm{eff}} = H(\mathcal A)/H(\mathcal K)\)};
			\node[label, right] at (3.0, 0.0) {\(K_{\mathrm{eff}} \propto m\)};
			\node[label, right] at (2.1, -1.6) {\ar{\(K_{\mathrm{eff}}\) = كفاءة الإصلاح}};
			\node[label, below] at (0.0, -2.5) {\ar{\(K_{\mathrm{eff}}\) = تنسيق الإجهاد}};
			\node[label, left] at (-2.1, -1.6) {\ar{\(K_{\mathrm{eff}}\) = شرط الاستقرار}};
			\node[label, left] at (-3.0, 0.0) {\(K_{\mathrm{eff}} \propto p_n\)};
			\node[label, left] at (-2.1, 1.6) {\ar{\(K_{\mathrm{eff}}\) = حجم القاموس}};
			\node[label, above] at (0.0, 2.5) {\ar{\(K_{\mathrm{eff}}\) = المعامل المتري}};
			
			% Внешняя рамка и заголовок
			\draw[gray!30, thick, rounded corners] (-6.9,-5.5) rectangle (6.9,5.8);
			\node[font=\small\bfseries, anchor=north west] at (-5.0,6.5) 
			{\ar{قانون التنسيق العالمي عبر جميع المجالات}};
			
		\end{tikzpicture}
		\caption{\ar{نفس الأس \(\beta = 2/3\) يتحكم في نظرية المعلومات، فيزياء الجسيمات، علم الأحياء، الفيزياء الأرضية، الديناميكا المائية، نظرية الأعداد، ميكانيكا الكم، وعلم الكونيات. في كل مجال، يأخذ \(K_{\mathrm{eff}}\) تفسيراً محدداً، لكن الشكل الدالي \(\varepsilon = \kappa_c \alpha K_{\mathrm{eff}}^{-2/3}\) يبقى ثابتاً. هذه العالمية هي علامة حقل تنسيق واحد أساسي \(\Psi_{MN}\).}}
		\label{fig:unified-beta}
	\end{figure}
	
	\section{\ar{تنبؤات قابلة للاختبار}}
	\label{sec:predictions}
	
	\ar{استناداً إلى القانون \(\eps = \kappac\,\alpha\,\Keff^{-2/3}\) مع \(\kappac = 1/3\)، نقترح التجارب التالية:}
	
	\begin{enumerate}
		\item \textbf{\ar{التحفيز الإنزيمي:}} \ar{لمجموعة من طفرات إنزيم، قِس \(\Keff\) (مثلاً من عدد الحالات التوافقية) والكفاءة التحفيزية \(k_{\text{cat}}/K_M\). يجب أن يكون ميل مخطط \(\log(k_{\text{cat}}/K_M)\) مقابل \(\log \Keff\) هو \(2/3\). التنبؤات البيولوجية الأكثر تفصيلاً موجودة في الملحقين D و E.}
		\item \textbf{\ar{الترميز العصبي:}} \ar{في مجموعة عصبية، قدّر \(\Keff\) من المعلومات المتبادلة بين المنبه والاستجابة مقسومة على إنتروبيا الاستجابة. يجب أن تتحجّم عتبة التمييز (معدل الخطأ) كـ \(\Keff^{-2/3}\).}
		\item \textbf{\ar{الأسواق المالية:}} \ar{باستخدام بيانات عالية التردد، احسب \(\Keff\) من سيولة دفتر الطلبات وفارق سعر العرض والطلب. من المتوقع أن تتبع التقلبية (الانحراف المعياري للعوائد) \(\sigma \propto \Keff^{-2/3}\).}
		\item \textbf{\ar{البنية الكونية:}} \ar{يجب أن يرتبط تموج الكثافة الفعال \(\sigma_8\) بكفاءة تنسيق المادة المظلمة، المقدرة من عدد المناطق المستقلة في الكون المبكر. وهذا يوفر فحص اتساق مع مشاهدات CMB (انظر النص الرئيسي والملحق G).}
		\item \textbf{\ar{القرص الدوار فائق التوصيل (الملحق Ehrenfest):}} \ar{بالنسبة لقرص مصنوع من موصل فائق ذي درجة حرارة عالية، يجب أن تختلف السرعة الزاوية الحرجة التي يتفكك عندها بين الحالة العادية والحالة فائقة التوصيل، لأن كفاءة تنسيق المادة \(K_{\mathrm{eff,mat}}\) تتغير. هذا يوفر اختباراً معملياً مباشراً للاعتماد المادي للهندسة.}
	\end{enumerate}
	
	\section{\ar{القيود والتحذيرات}}
	\label{sec:limitations}
	
	\ar{يجب توخي الحذر عند تطبيق قانون التحجيم العالمي:}
	\begin{itemize}
		\item \ar{تعريف \(\Keff\) يتطلب تحديداً واضحاً للقاموس وفضاء الأفعال؛ في العديد من الأنظمة، هذا ليس بالأمر التافه.}
		\item \ar{العامل المقدم \(\kappac = 1/3\) دقيق للبنية الثلاثية. في YUCT، \(\alpha\) هو ثابت ثابت مشتق في الملحق Y، القسم Y.5.2؛ وليس معاملاً حراً. ومع ذلك، عملياً، عندما يُقدّر \(\Keff\) تقريبياً فقط، قد تظهر القيمة الفعالة لـ \(\alpha\) المستنتجة من البيانات متغيرة. هذا يشير إلى الحاجة إلى تقديرات أكثر دقة لـ \(\Keff\)، وليس إلى تباين أساسي في \(\alpha\).}
		\item \ar{الأس \(2/3\) مشتق بدقة من البعد الثلاثي للفضاء (النظرية 2) وشرط الاتصال (الملحق Y)، لكن الإسقاط من \(\Keff\) إلى الكميات القابلة للرصد قد يتطلب افتراضات إضافية.}
		\item \ar{تتضمن الأمثلة في الجدول~\ref{tab:data} أنظمة حيث الخطأ ليس الكمية محل الاهتمام المباشر (مثل CMB)؛ لذا يلزم إسقاط دقيق.}
	\end{itemize}
	
	\ar{ومع ذلك، فإن تقارب البيانات من أكثر من 50 مرتبة مقدارية يشير بقوة إلى وجود ظاهرة عالمية حقيقية.}
	
	\section{\ar{الاستنتاج}}
	\ar{قدمنا صياغة منقحة ومتسقة رياضياً لقانون تحجيم الأخطاء العالمي ضمن YUCT. الأس \(\beta = 2/3\) دقيق ويتبع من البعد الثلاثي لفضاء التنسيق (النظرية 2 في النص الرئيسي) وشرط اتصال التسلسل الهرمي (الملحق Y)، بينما يعكس العامل المقدم \(\kappac = 1/3\) البنية الثلاثية. القيمة التجريبية \(0.67 \pm 0.03\) تؤكد هذا التوقع. يوحد القانون مجموعة واسعة من قوانين القوى في الأنظمة المعقدة ويقدم تنبؤات قابلة للاختبار. يجب أن تركز الأعمال المستقبلية على القياسات الدقيقة لـ \(\Keff\) في مجالات مختلفة وتوسيع الاشتقاق النظري ليشمل العامل المقدم \(\alpha\).}
	
	\paragraph*{\ar{توفر البيانات}}
	\ar{جميع تجميعات البيانات ونصوص التحليل متاحة على} \url{https://github.com/Alexey-Yakushev-YUCT/YPSDC}.
	
	\paragraph*{\ar{تاريخ الإصدار}}
	\begin{itemize}
		\item v1.0: \ar{نسخة أولية بتعاريف غير متسقة.}
		\item v2.0: \ar{تصحيح التعاريف، إضافة أمثلة تحقق.}
		\item v2.2: \ar{توسيع إلى المقاييس الذرية والبيولوجية.}
		\item v3.0: \ar{مراجعة الاشتقاق باستخدام النظرية 2، استبدال القيمة التجريبية \(0.67\) بالأس الدقيق \(2/3\)، توضيح القيود، إضافة مراجع للملاحق P، Y، Ferra1.2.}
		\item v3.1: \ar{توضيحات إضافية: إضافة اشتقاق صريح للعامل 2 من الملحق Y، حذف مثال تحلل النيوترون المضلل، إعادة صياغة ارتباط نموذج إيزينغ كتأكيد مستقل، إضافة مراجع للملاحق D، E، F للتشغيل، تصحيح مناقشة \(\alpha\) لتعكس وضعها كثابت ثابت.}
		\item v3.2: \ar{إضافة قسم التنسيق العالمي عبر المجالات، بما في ذلك الهياكل الحلزونية والأعداد الأولية؛ ربط بالملحقين Ehrenfest و X؛ إضافة مراجع ببليوغرافية ذات صلة.}
		\item v3.3: \ar{إضافة قسم فرعي توليفي يوحد جميع المجالات بشكل صريح، وشكل ملخص بصري (الشكل~\ref{fig:unified-beta}) يظهر الدور المركزي لـ \(\beta = 2/3\) في نظرية المعلومات، فيزياء الجسيمات، علم الأحياء، الفيزياء الأرضية، الديناميكا المائية، نظرية الأعداد، ميكانيكا الكم، وعلم الكونيات.}
	\end{itemize}
	
	\begin{thebibliography}{99}
		\bibitem{Yakushev2026}
		A. V. Yakushev, \emph{Yakushev Unified Coordination Theory (YUCT)}, Zenodo, \url{https://doi.org/10.5281/zenodo.18444599} (2026).
		\bibitem{AppendixC1}
		A. V. Yakushev, \ar{الملحق C1: التحقق التجريبي من قانون التحجيم العالمي في الروابط الكيميائية (HF–HI)، في} \emph{YUCT}، Zenodo (2026).
		\bibitem{AppendixD}
		A. V. Yakushev, \ar{الملحق D: التنبؤات البيولوجية لـ YUCT – فرضيات قابلة للاختبار، في} \emph{YUCT}، Zenodo (2026).
		\bibitem{AppendixE}
		A. V. Yakushev, \ar{الملحق E: علم الوراثة التنسيقي – مبادئ YUCT في البيولوجيا الجزيئية، في} \emph{YUCT}، Zenodo (2026).
		\bibitem{AppendixF}
		A. V. Yakushev, \ar{الملحق F: تطبيق YUCT على الاقتصاد، في} \emph{YUCT}، Zenodo (2026).
		\bibitem{AppendixG}
		A. V. Yakushev, \ar{الملحق G: نظرية التنسيق الموحدة لياكوشيف – حل أساسي لمشكلة الجاذبية الكمية، في} \emph{YUCT}، Zenodo (2026).
		\bibitem{AppendixP}
		A. V. Yakushev, \ar{الملحق P: الاعتماد الحراري للجاذبية، في} \emph{YUCT}، Zenodo (2026).
		\bibitem{AppendixY}
		A. V. Yakushev, \ar{الملحق Y: الأسس الرياضية لـ YUCT – اشتقاق من المبادئ الأولى، في} \emph{YUCT}، Zenodo (2026).
		\bibitem{AppendixFerra}
		A. V. Yakushev, \ar{الملحق Ferra1.2: الثوابت التنسيقية العالمية للأطوار المرتبة وغير المرتبة، في} \emph{YUCT}، Zenodo (2026).
		\bibitem{AppendixPrimeN}
		A. V. Yakushev, \ar{الملحق PrimeN: سلم تنسيق الأعداد الأولية لـ YUCT، في} \emph{YUCT}، Zenodo (2026).
		\bibitem{AppendixX}
		A. V. Yakushev, \ar{الملحق X: نظرية شانون المعممة، قانون الخطأ الثابت، وحد بل المعتمد على \(K_{\mathrm{eff}}\)، في} \emph{YUCT}، Zenodo (2026).
		\bibitem{AppendixEhrenfest}
		A. V. Yakushev, \ar{الملحق Ehrenfest: التفسير التنسيقي لمفارقة القرص الدوار وظهور الهندسة غير الإقليدية، في} \emph{YUCT}، Zenodo (2026).
		\bibitem{El-Showk2014}
		S. El-Showk et al., \emph{Phys. Rev. D} \textbf{90}, 105001 (2014).
		\bibitem{ammonites}
		D. K. Jacobs, Morphology and the evolution of ammonite shell coiling, \emph{Paleobiology} \textbf{16}, 283 (1990).
		\bibitem{watsoncrick}
		J. D. Watson and F. H. C. Crick, Molecular structure of nucleic acids, \emph{Nature} \textbf{171}, 737 (1953).
		\bibitem{pauling}
		L. Pauling and R. B. Corey, Configurations of polypeptide chains with favored orientations around single bonds, \emph{Proc. Natl. Acad. Sci. USA} \textbf{37}, 729 (1951).
		\bibitem{collagen}
		G. N. Ramachandran and G. Kartha, Structure of collagen, \emph{Nature} \textbf{176}, 593 (1955).
		\bibitem{phyllotaxis}
		P. H. Rivier and A. M. Turing, The chemical basis of morphogenesis, \emph{Phil. Trans. R. Soc. B} \textbf{237}, 37 (1952).
		\bibitem{cholesteric}
		P. G. de Gennes and J. Prost, \emph{The Physics of Liquid Crystals}, Oxford University Press, 2nd ed., 1993.
		\bibitem{helical-mof}
		B. Moulton and M. J. Zaworotko, From molecules to crystal engineering: supramolecular isomerism and polymorphism in coordination polymers, \emph{Chem. Rev.} \textbf{101}, 1629 (2001).
		\bibitem{ocean-eddies}
		A. R. Robinson (ed.), \emph{Eddies in Marine Science}, Springer, 1983.
		\bibitem{hurricanes}
		K. Emanuel, The physics of tropical cyclones, \emph{Annu. Rev. Fluid Mech.} \textbf{35}, 1 (2003).
		\bibitem{tornadoes}
		R. Davies-Jones, A review of supercell and tornado dynamics, \emph{Atmos. Res.} \textbf{158--159}, 274 (2015).
		\bibitem{nuclear-tracks}
		R. L. Fleischer, P. B. Price, and R. M. Walker, \emph{Nuclear Tracks in Solids: Principles and Applications}, University of California Press, 1975.
		\bibitem{binney-tremaine}
		J. Binney and S. Tremaine, \emph{Galactic Dynamics}, Princeton University Press, 2nd ed., 2008.
		\bibitem{ppdisk}
		S. M. Andrews and D. J. Wilner, The structure of protoplanetary disks, \emph{Annu. Rev. Astron. Astrophys.} \textbf{55}, 471 (2017).
		\bibitem{beck}
		R. Beck, Magnetic fields in galaxies, \emph{Space Sci. Rev.} \textbf{99}, 243 (2001).
		\bibitem{Tegmark2014}
		Max Tegmark, \emph{Our Mathematical Universe: My Quest for the Ultimate Nature of Reality}, Knopf, 2014.
	\end{thebibliography}
	
\end{document}